% MT014 / explanation — canonical editable content.
生成式和判别式都要给新输入分类，区别在于先学什么。生成式分别描述“各类通常产生怎样的数据”，再比较类别；判别式直接学“输入怎样对应到标签”，不必完整描述每一类的数据分布。本题要求选\textbf{不正确}的说法。



\begin{itemize}

\item \textbf{A错误，选。} 假设原来根据长度分类，现在新增重量。模型至少要知道每类重量的大致分布，否则新的一列数字无法提供类别证据。NB中原有长度参数可以保留，但重量参数仍需学习；若联合模型还描述特征关系，相关参数也要更新。

\item \textbf{B正确，不选。} 生成式学习\(p(x\mid c)\)及\(P(c)\)，预测时比较\(p(x\mid c)P(c)\)。前者说明新样本像不像该类，后者考虑该类本来有多常见。

\item \textbf{C按本课监督分类设定错误，选。} 训练数据要给出正确类别，才能知道当前判别结果对不对。无标签预训练、半监督方法属于额外流程，不能拿来改变此题的标准监督设定。

\item \textbf{D错误，选。} 分布必须匹配特征含义。身高这类连续测量可以尝试高斯分布；“词是否出现”是0/1，更适合Bernoulli模型；词次数可用Multinomial模型。不存在高斯在所有场合都更好的结论。

\item \textbf{E按课程分类正确，不选。} 题中Bayes分类器走“分布与先验\ensuremath{\to}后验”的路线；SVM、LR及这些树集成方法直接学习分类规则或类别概率。

\end{itemize}


\textbf{出处：} 2021A期中 Q1，原卷第2页（5分）


\textbf{答案依据：} 2021A配套解答 / 2021A配套解答副本 Q1，随卷答案ACD（题目要求选不正确）。
